\begin{table}[htbp]
\centering
\footnotesize
\caption{Uncertainty Around the Architecture-vs-Information Decomposition}
\label{tab:disentangling_uncertainty}
\begin{tabular}{lccc}
\toprule
Contrast & Estimate & 95\% CI & Excludes 0 \\
\midrule
Architecture average & 477.8 & [221.3, 676.1] & Yes \\
Architecture at Parsimonious & 384.8 & [180.1, 538.9] & Yes \\
Architecture at BoP & 570.9 & [244.7, 838.1] & Yes \\
Information average & 78.1 & [-3.1, 149.4] & No \\
Information at MS-VAR & -14.9 & [-54.9, 21.6] & No \\
Information at XGBoost & 171.2 & [-5.3, 315.5] & No \\
Interaction (DiD) & 186.1 & [-9.4, 345.0] & No \\
\bottomrule
\end{tabular}
\begin{tablenotes}
\footnotesize
\item \textit{Notes:} Positive architecture effects mean XGBoost has higher RMSE than MS-VAR, so larger values favour the regime-switching architecture. Positive information effects mean the BoP specification raises RMSE relative to the parsimonious specification under the paper's sign convention. Confidence intervals use a stationary block bootstrap over rolling-origin dates (2023:01--2025:03; $n=27$, 1000 replications, block length 3).
\end{tablenotes}
\end{table}
